In the realm of physics, symmetries are fundamental for unraveling the conservation laws that govern physical systems. However, it is unfortunate that these symmetries can be disrupted.

There are three distinct pathways through which a symmetry, whether continuous or discrete, can be compromised during the investigation of a system. This phenomenon is referred to as "the breaking of a symmetry."

\begin{enumerate}
    \item \textbf{Explicit Symmetry Breaking}:
          \begin{itemize}
              \item \textbf{Description}: Explicit symmetry breaking occurs when external influences or explicitly added terms in a physical system break the symmetry that the system would otherwise possess.
              \item \textbf{Example}: An external magnetic field applied to a material can explicitly break the rotational symmetry of the material's electrons.

                    Or, in a more abstract example, but one that resembles the problem we will discuss: a free particle on the whole real axis obviously possesses parity symmetry and translational symmetry. The introduction of a potential lacking parity unequivocally disrupts both symmetries.
          \end{itemize}

    \item \textbf{Spontaneous Symmetry Breaking}:
          \begin{itemize}
              \item \textbf{Description}: Spontaneous symmetry breaking occurs when the laws governing a physical system are symmetric, \textit{e.g. the Lagrangian}, but the system's ground state (the lowest energy state) does not present the same symmetry.

              \item \textbf{Example}: In a ferromagnetic material \textit{e.g. iron}, the individual magnetic moments of its atoms tend to align in a specific direction. At high temperatures, these magnetic moments are randomly oriented, and the material has no net magnetization. In this state, there is a symmetric orientation of magnetic moments.

                    However, as the temperature decreases, there is a critical temperature called the Curie temperature. Below this temperature, the material spontaneously undergoes a quantum mechanical phase transition, and the magnetic moments align in the same direction, resulting in a net magnetization. The symmetry is spontaneously broken as the material acquires a macroscopic magnetization in a specific direction.

                    This phenomenon is a classic example of spontaneous symmetry breaking in classical physics and is responsible for the behavior of permanent magnets and the Earth's magnetic field.

                    Or, in a more abstract example, but one that resembles the problem we will discuss: a particle in a symmetric double well potential. Even though the system possesses parity symmetry, in Quantum Mechanics, the particle spontaneously chooses one well over the other, breaking the symmetry.
          \end{itemize}

    \item \textbf{Anomalous Symmetry Breaking}:
          \begin{itemize}
              \item \textbf{Description}: Anomalous symmetry breaking is a phenomenon where the process of quantization of a system, \textit{i.e. to solve it using quantum mechanics}, break symmetries that would be preserved in a classical theory.
              \item \textbf{Example}: In quantum field theory, anomalies can result in the violation of classically conserved currents, such as the chiral anomaly in the context of the strong force, which breaks specific symmetries.

                    In our case, the scale invariance of the problems will be broken when the Hamiltonian operator of the problem, originally lacking physical significance, is elevated to the status of a physical Observable.
          \end{itemize}
\end{enumerate}

\noindent
The first two symmetry breaking mechanisms are well known to physicists since the apply to classical as well as quantum systems.
Rather, our discussion will focus on the subject of anomalous symmetry breaking.

Typically, this mechanism is encountered within Quantum Field Theory, such as the chiral anomaly in processes like $\pi^0 \rightarrow 2 \gamma$. However, there are two examples from ordinary quantum mechanics: the two-dimensional $\delta$-potential and the $-1/x^2$ potential. We will study the latter, but as mentioned before, we will establish a connection between them at the end of the essay in Section~\ref{section:delta_potential}, demonstrating that they are not entirely distinct problems.